\chapter{Conclusão}
\label{cap:conclusao}

Este trabalho apresenta uma sumarização dos resultados necessários da geometria Riemanniana para enunciar e demonstrar propriedades acerca de geodésicas em variedades Riemannianas, incluindo o Teorema de \nameref{teo:hopf-rinow}. Desenvolvemos a noção de variedade Riemanniana e munimos toda variedade Riemanniana da conexão de Levi-Civita, a qual é simétrica e compatível com a métrica. Utilizamos isso para desenvolver a noção de geodésica e apresentar alguns resultados, incluindo o Teorema de \nameref{teo:hopf-rinow}. Esse teorema apresenta algumas equivalências em variedades Riemannianas conexas, dentre as quais temos que o mapa exponencial está bem definido em um ponto, para todo vetor no espaço tangente, a variedade é completa na métrica, no sentido de que toda sequência de Cauchy converge e a variedade é geodesicamente completa, isto é, quaisquer dois pontos podem ser ligados por uma geodésica minimizante a qual pode ser extendida indefinidamente. Essa última equivalência representa a passagem do local a global do fato de que, em qualquer variedade Riemanniana, existem geodésicas minimizantes localmente. Fica evidente a pertinência do estudo das geodésicas, além de resultados obtidos a partir dele.
